%%%%%%%%%%%%%%%%%%%%%%%%%%%%%%%%%%%%%%%%%%%%%%%%%%%%%%%%%%%%%%%%%%%%%%%%%%%%
%                                                                          %
%       Numerical differentiation in system theory. Observer design        %
%                                                                          %
%       Edited on August 8, 1993                                           %
%%%%%%%%%%%%%%%%%%%%%%%%%%%%%%%%%%%%%%%%%%%%%%%%%%%%%%%%%%%%%%%%%%%%%%%%%%%%

%%%%%%%%%%%%%%%%%%%%%%%%%%%%%%%%%%%%%%%%%%%%%%%%%%%%%%%%%%%%%%%%%%%%%%%%%%%%
%   Comment out and adjust the following lines when using Double Columns   %
%%%%%%%%%%%%%%%%%%%%%%%%%%%%%%%%%%%%%%%%%%%%%%%%%%%%%%%%%%%%%%%%%%%%%%%%%%%%
% \hoffset = -10 true mm
% \voffset = -5 true mm
% \hsize = 185 true mm
% \vsize = 235 true mm
% \newdimen\separatorwidth
% \separatorwidth = 5 true mm
% \tolerance = 1000
% \emergencystretch = 25pt
%%%%%%%%%%%%%%%%%%%%%%%%%%%%%%%%%%%%%%%%%%%%%%%%%%%%%%%%%%%%%%%%%%%%%%%%%%%%

\input astm

\nopagenumbers
\footline = {\ifnum \pageno = 1 \else \hss \eightpoint -- \folio\ -- \hss \fi} 

\centerline{\seventeenbf Numerical differentiation in system theory.}
\centerline{\seventeenbf Observer design\footnote\fnast{\sevenpoint
This is the very first draft of a work in preparation by S. Diop, J.~W.
Grizzle and P.~E. Moraal.}}

\bigskip\noindent
{\twelvebf Abstract.} To be re-written.  I moved part of it to
the Introduction (J.G.).

Numerical differentiation is well known as an ill-conditioned problem in the
sense that small perturbations on the function to be differentiated may induce
large ones on the derivations.
We illustrate the fact that some specific methods are able to overcome this
difficulty, turning numerical differentiation techniques into a
challenger to existing  observer design methods.

\medskip\noindent
{\twelvebf Keywords.}
Observer design; Numerical differentiation; System theory. 

\bigskip\noindent
{\twelvebf Introduction.}  Observer design is a fundamental problem in system theory and 
control practice.  Indeed, many feedback  design techniques, especially for nonlinear
system models,  start with the assumption that the full state is available for
feedback and that the {\it user} will design separately a suitable observer
for the actual implementation.  This in turn has motivated a great deal of
research into the observer design problem, for which 
we will not attempt to give a survey of available results.   

The specific problem addressed herein involves the situation where
one seeks to ``estimate'' the states of a continuous-time model from
observations collected at discrete instances in time.  To date, this problem has
been addressed only in the context of sampled data systems \cite{GM, MG}; the
procedure has been to first compute
a discrete-time model from the continuous-time model, and
then to design an observer on the basis of the discretized system model.
Of course, the discretization process sometimes can be quite non-trivial, from
a numerical or computational point of view; also, it is usually quite important 
to assume that the observed data has been collected at regular instances in time,
and thus dealing
with problems where event driven sampling is necessary can be difficult
\cite{??}.  

This paper comes out of the observation that
on one hand, the numerical analysis literature contains a
substantial body of results on numerical differentiation, while on
the other hand, one of the very basic problems in system theory turns
out to be that of obtaining the time derivatives of a continuous system
variable, which is usually known only through the differential equations of
the system, and time samples of this variable.
We present a first attempt in applying results on numerical differentiation
to some system theoretic problems; in this paper,
we focus on the observer design problem, mainly by
working out some specific examples.

To see how numerical differentiation and observer design may tie together, 
let us consider the 
following simple observer design problem.
Let a system be given as
$$
\left\{
\eqalign{
\dot x_1 & = x_2\cr
\dot x_2 & = -ux_1^3\cr
y & = x_1\cr}
\right.
\eqno(1)
$$
and let the question be to build an observer for $(1)$.
It is clear that the variable $x_1$ need not be computed since it 
is directly given by the measurement: $ x_1 = y$.
From the equations of $(1)$, we see that $x_2 = \dot y$.
If, as is usually the case, we have no access to 
$\dot y$ then we need to {\it compute} $\dot y$ from the measurement $y$.
This is where numerical differentiation comes into play:
we explore the feasibility of the solution which 
consists in directly {\it numerically
differentiating} $y$ in order to obtain $x_2$.

If $y$ is known only through its time samples then numerical 
differentiation yields an approximation of
$\dot y$ at the current sampling time $t_k$ by some number
$\hat {\dot y}(t_k)$.
An algorithm for the computation of $\hat {\dot y}(t_k)$ from the
samples of $y$, and {\it estimates of the error bounds} should thus be worked out.
Usually, not only $\dot y$, but, also, finitely many higher
derivatives of $y$ should be approximated at time $t_k$ in order to be
able to compute all of the desired state components.

The intent of this simple example is to 
illustrate our main objectives in applying 
numerical differentiation techniques to the observer design problem.
If efficient numerical 
differentiation algorithms are available, then they should be
applicable to many other system theoretic problems requiring the
computation of the derivatives of continuous variables which are
known only through their time samples and the differential equations they
satisfy.
The inversion problem and some control problem such as
regulation are among other potential applications we have in mind.

Our investigation led us directly to interpolation theory, and more
generally, to approximation theory.
We thus noted that interpolation theory has already been used in system
theory literature, mainly for constant linear system realization
purposes, in
\cite{youla1967,antoulas1986,antoulas1988,antoulas1989a,antoulas1990a,anderson1990b}.
See also \cite{hofmann1977}.  {\bf Is this really related??}

As is well known, there is a considerable number of different ways for
approximating the value of an unknown function at some point.
One commonly chooses among polynomial, rational, periodic, or exponential
functions
as basis elements.
(See also methods stemming from the Shannon sampling theorem
\cite{butzer1971,butzer1977,splettstosser1979,splettstosser1983,butzer1983a,butzer1983b,ries1983,schussler1983}.)
The optimality of any approximation process is, of course, relative to such
assumed base functions.
Approximation methods based on polynomials are the
simplest ones, but spline approximations seem to have some interesting
fundamental properties.


{\it Tornambee: talk about his work.}
As far as we know, system theorists used to invoke asymptotic 
observer design techniques to solve this problem.
While this approach is very successful for constant linear systems, the
nonlinear case seems restrictive, and even obstructive.


\bigskip\noindent
{\twelvebf 1. An observability property.}
Given a system described by the state equations
$$
\left\{
\eqalign{
\dot x_i & = g_i(x, u) \quad (1 \le i \le n)\cr
y_j & = h_j(x, u) \quad (1 \le j \le p)\cr}
\right.
\eqno(2)
$$
we will adopt the following specific observability property (see
\cite{chung1990,grizzle1990b,grizzle1990c,grizzle1990a,moraal1992,song1992,moraal1993b,moraal1993a}):
system $(2)$ is {\bf observable} if there is an integer $N$ such
that the map $H\left(u, \dot u, \cdots, u^{(N-2)}, \cdot \right)$ of
the state space into some space of output values, defined by $$
H\left(u, \dot u, \cdots, u^{(N-1)}, x \right) = \pmatrix{y \cr \dot
y \cr \vdots \cr y^{(N-1)} \cr}, \eqno(3) $$
is {\it injective} for any fixed {\it universal input} \cite{Sussmann??}.
We shall define later what these universal inputs are.
Note that it is easier at this stage to assume that there is actually
no input in $(2)$, or that all inputs are universal in the sense that
the above map is injective for any input. 

For an observable system $(2)$, we then may write
$$x = L\left(u, \dot u, \cdots, u^{(N-2)}, y, \dot y, \cdots, y^{(N-1)} 
\right).$$
The existence of the map $L$ is guaranteed by our definition of 
observability, but an explicit expression of $L$ may not be easy to 
obtain so that numerical equation solving methods such as Newton's algorithm
may be required in order to obtain $L$.  This is not the focus of this
paper.
When $(2)$ is a rational system in the sense that the $g_i$'s and the 
$h_j$'s are rational functions of their arguments, then the type of 
observability we just defined corresponds to {\it rational 
observability} of \cite{diop1991d,diop1992b}, and then the explicit
expression of $L$ may be  constructively derived through the use of cetain
differential
algebraic techniques.

\noindent {\bf Note:} In the literature, it is often tacitly assumed that $N$ should
be taken  as $2n + 1$; we shall examine this situation which, {\it a
priori}, seems connected to the works \cite{aeyels1981a,aeyels1981b}.

For these observable systems, the observer design problem may thus be 
seen as a problem of numerical differentiation: once estimates of the
derivatives
of y and u can be detrmined from the available measurements, then x can be
determined from $L$.

\bigskip\noindent
{\twelvebf 2. On numerical differentiation.}
The basic problem we are facing may be outlined as following.

\proclaim Problem P.
A real valued function $f$ of one real variable
(called the  {\it time,} $t$) defined on an interval $\BI$ is known
only through the differential equations it satisfies and through the
values it  takes at some instants $t_0, t_1, \cdots, t_k \in \BI$,
with $t_0 <$ $t_1 < \cdots <$ $t_k$.  The question is to provide
estimates of the first  $q$ derivatives $f^{(1)}(t_k)$,
$f^{(2)}(t_k), \cdots$, $f^{(q)}(t_k)$ of $f$ at the time $t_k$.
The difficulty of the problem may be increased by the presence of noise in the samples
of $f$, in which case, the approximation technique should provide
{\it smoothing}, and even the values of $f$ would need to be estimated.

\noindent {\bf Remarks:}

\begin{itemize}

\item
The more general problem we should end up facing to in system
theory is actually a multivariable one in that $f$ is a vector function
instead of a scalar one as in the latter statement.

\item
If $r$ is the order of the differential equations defining $f$ then only
at most the first $q = r - 1$ derivatives of $f$ need be estimated, the
succeeding derivatives being computable from the previous ones
by means of these differential equations defining $f$.

\item
Problem~P is a particular case of the general problem of approximation of
functions on the basis of partial knowledge on these functions.
Its seems that there is not yet a complete solution with
expressions of the error bounds in the literature.
One partial solution consists in taking the derivatives of the
{\it interpolant} $\hat f$ for those of $f$.
In this standpoint there is a substantial amount of literature.
Another possible direction consists in trying to overcome
the ill-conditionedness of the numerical differentiation problem by using
some regularization techniques.
The only known works are the one by Cullum~\cite{cullum1971},
and those in the Russian literature (cited in the latter paper) to which we
do not yet have access.
\end{itemize}

\medskip\noindent
{\elevenbf 2.1. Generalities on approximation theory.}
The basic problem that approximation theory addresses is the following:
Assume that the unknown object to be approximated is an element
$x$ of a given metric space $\BE$, and find a {\it best estimate}
$\hat x$ of $x$ in an (also given) subspace $\BF$ of $\BE$ subject to
$$
{\rm d}(x,\hat x) = {\rm d}(x,\BF) = \min_{y \in \BF}{\rm d}(x,y),
$$
where {\rm d} denotes the distance on $\BE$.
The word {\it best} is of course relative to the distance used in
$\BE$ and to the choice of $\BF$.
For example, for a fixed distance, the larger $\BF$ is, the finer the
approximation will be.

We know that this problem has at least one solution if the subspace
$\BF$ is compact, or if $\BE$ is a normed vector
space and $\BF$ is a linear subspace of finite dimension, or if $\BE$
is a Hilbert space and $\BF$ is a closed subspace.

Inour case, the object to be approximated is the set of derivatives $f^{(1)}$,
$f^{(2)}, \cdots$, $f^{(q)}$ of a real valued function $f$ defined
on, say a compact interval $\BI$ of $\setR$, so that the space $\BE$
may be taken as the set of real valued functions which are $q-1$ times
differentiable with $f^{(q-1)}$ absolutely continuous and $f^{(q)}$ square
integrable on $\BI$, or which are infinitely differentiable, or, even, which
are analytic on $\BI$.
$\BE$ is thus endowed with its usual $\setR$-algebra
structure and may be convertible into a real normed vector space by one of
the following norms:
$$
\eqalign{
\| g \|_p & = \left( \int_\BI |g(t)|^pdt \right)^{1/p}, \, g \in \BE, p
\in \setR, \, p \ge 1, \cr
\| g \|_\infty & = \sup_{t \in \BI} |g(t)|, \, g \in \BE, \cr
\| g \|_{(q)} & = \| g \|_2 + \| g^{(1)} \|_2 + \cdots +
\| g^{(q)}\|_2, \, g \in \BE, \cr
\hbox{\rm etc.} & \cr}
$$
The norm $\| \, \|_2$ is the one of a Hilbert space
structure on $\BE$, the corresponding inner product being
$$
\left( g, h \right) = \int_I g(t)h(t)dt, \, g, h \in \BE.
$$
The norm $\| \, \|_\infty$ is usually called the {\it infinity} or {\it
uniform} norm.

For example, given that $\BF$ is a finite dimensional subspace of
$\BE$, the {\it least square} approximation process assumes the norm
$\| \, \|_2$, and results in the existence and uniqueness of the best
estimate of any given $f$ in $\BE$.

The subset $\BF$ of $\BE$ is usually taken as a subspace of one of
the following linear subspaces:
\item{$\bullet$}
$\setR\left[ t \right]$, the subspace of polynomial functions,
\item{$\bullet$}
$\setR\left( t \right)$, the subspace of rational functions,
\item{$\bullet$}
$\setR\left[ \cos t, \sin t \right]$, the subspace of $\setR$-linear
combinations of $\cos(it)$ and $\sin(it)$, $i \in \setN$,
\item{$\bullet$}
etc.

Assume linear uniform approximation theory,
and $\BF$ finite dimensional, with dimension $N'+1$.
We then know that, given any $f$ in $\BE$, there is at least one
best estimate $\hat f$ of $f$.

A base $(\zeta_0$, $\zeta_1, \ldots$, $\zeta_{N'})$ of $\BF$ is said to
satisfy the Haar condition if any real linear combination
$$
c_1\zeta_1(t) + \cdots + c_N\zeta_{N'}(t)
$$
vanishes at less than $N'+1$ distinct points in $\BI$.

\proclaim Theorem 1 {\rm (Haar), \cite{delavalleepoussin1970}}.
A necessary and sufficient condition for the best uniform estimate to be
unique is that the base $(\zeta_0$, $\zeta_1, \ldots$, $\zeta_{N'})$ of $\BF$
satisfies the Haar condition.

\medskip\noindent
{\elevenbf 2.2. Differentiating the interpolant.}
Since the interpolation problem has been given a large number of elegant
and practical solutions, the easiest way to approach Problem~P is to take
the derivatives of the interpolant $\hat f$ for those of $f$.
Let us put the following remark at this stage even if we have not yet
introduced the material which makes it clear.

There might be, due the reality of the
physical problem, no freedom in the choice of the sampling instants $t_i$,
but the error bound is reasonably expected to depend on both the number of
samples and  their distribution on the time interval $\BI$ as apparent from
the Tchebychev theorem.
Our study should address this problem, as well, but
for now, we shall suppose that the current time is $t_N$, later on, we shall
discuss how we should deal with the succeeding sampling instants $t_{N+1},
\ldots$ as time evolves.

Let $\BF_{N'}$ be a linear subspace of $\BE$ of finite dimension $N'+1$, and
with base $(\zeta_0$, $\zeta_1, \ldots$, $\zeta_{N'})$.
An element
$$
\hat \zeta_{N'}(t) = c_0\zeta_0(t) + \cdots + c_{N'}\zeta_{N'}(t)
$$
of $\BF_{N'}$ {\it interpolates} $f$ if
$$
\hat \zeta_{N'}(t_i) = f(t_i) \, (0 \le i \le N),
$$
that is, if
$$
\sum_{j=1}^{N'}c_j\zeta_j(t_i) = f(t_i) \, (0 \le i \le N).
$$
Therefore,

\proclaim Theorem 2.
There is one and only one function in $\BF_{N'}$ interpolating $f$
if and only if the previous system in the $c_j$'s has one and only one
solution, that is, the existence and uniqueness in $\BF_{N'}$ of an
interpolant of $f$ is equivalent to $N' = N$ and the determinant of the
following matrix
$$
\pmatrix{\zeta_0(t_0)&\ldots&\zeta_N(t_0)\cr
         \vdots&&\vdots\cr
         \zeta_0(t_N)&\ldots&\zeta_N(t_N)\cr}
$$
is nonsingular; in which case, we may use the well-known Cramer formulae to
compute the interpolant, and write it in the form
$$
\hat \zeta_N(t) = \sum_{i=1}^{N}f(t_i)s_i(t) \, (t \in \BI),
$$
with $s_i(t_j) = \delta_{ij}$, where $\delta_{ij}$ is the Kronecker symbol.

\noindent
Note that the Haar condition guarantees that the latter square matrix is
nonsingular.
Note also that this previous result shows that, when interpolating linearly a
function at $N$ points, the subspace of projection, $\BF$, should be of
dimension precisely $N$.

Among the well-established results on interpolation theory we have the
famous Lagrange one, the spline one, etc.

\smallskip\noindent
{\bf 2.2.1. Differentiating the Lagrange interpolant.}
It is clear that there are infinitely many polynomials of arbitrary degree
{\it interpolating} $f$, but, according to the previous general result,
there is one and only one of degree at most $N$ which interpolates $f$.
This unique polynomial $L_N$ is known as the Lagrange (or polynomial)
interpolant of $f$.

Assuming
$$
\ell_i(t) = \prod_{\scriptstyle j = 0 \atop \scriptstyle j \ne i}^N{t-t_j \over
t_i - t_j},
$$
we have
$$
L_N(t) = \sum_{i=0}^{N}f(t_i)\ell_i(t) \, (t \in \BI),
$$
with $\ell_i(t_j) = \delta_{ij}$.
This is not the most interesting form for the computation of the Lagrange
interpolant.
The Newton formulae lead to more efficient computations,
see~\cite{stoer1980,marchuk1982,nonweiler1984} for example.

If the only knowledge we have on $f$ is the set of its samples, then being
able to estimate the error bounds is hopeless. We have the following result:
If $f$ is $N$ times continuously differentiable in then for any $t$
in $\BI$, there is a point $\xi_t$ in $\left[ t_0, t_N \right]$ such that
$$
\varepsilon(t) = f(t) - \hat \zeta_{N}(t) = {1 \over
(N+1)!}\ell(t)f^{(N+1)}(\xi_t) \, (t \in \BI),
$$
where $\ell(t) = \prod_{i=0}^N(t-t_i)$.

This follows mainly from the Rolle theorem.
In general, given $t$, we do not know neither $\xi_t$ nor
$f^{(N+1)}(\xi_t)$.
This is the reason why we shall need to know a bound for the higher
derivatives of $f$:
$$
\| \varepsilon \|_\infty \le {\Delta^N \over (N+1)!}\| f^{(N+1)} \|_\infty,
$$
where $\Delta$ is $\max_{1 \le i \le N}(t_i -t_{i-1})$.
This bound depends both on the number $N+1$ of samples and on their
distribution.
If it is possible to choose freely the sampling instants, then we would be
able to decrease the latter error bound by sampling according to the
Tchebychev theorem, see \cite{nonweiler1984}.

Now, the error when taking $\hat \zeta_{N}^{(j)}(t)$ for $f^{(j)}(t)$ is
given by
$$
\varepsilon^{(j)}(t) = f^{(j)}(t) - \hat \zeta_{N}^{(j)}(t) = \sum_{\iota =
0}^j{1 \over (N+\iota+1)!}{j \choose
\iota}\ell^{(j-\iota)}(t)f^{(N+\iota+1)}(\xi_\iota) \, (t \in \BI),
$$
where the $\xi_\iota$'s are $t$ dependent and in $\left[ t_0, t_N
\right]$.
Therefore, we have
$$
\| \varepsilon^{(j)} \|_\infty \le c \, \Delta^{N+1-j} \max_{0 \le \iota \le
j}\| f^{(N+1+\iota)} \|_\infty,
$$
where $c$ is an absolute constant (it does not depend on neither $f$ nor
the samples), see \cite{schaback1992}.

\proclaim Remark.
Though simple and elegant, the Lagrangian interpolation suffers from the
{\rm Runge phenomenon} which is the fact that small change to the data
points in the middle of the data set may produce a very large excursion in
the end points, $t_0, t_N$, see \S 3.2.3 of \cite{nonweiler1984}.

According to the latter book,
Lagrangian interpolation is thus not recommended for large set of equispaced
data.
Nevertheless, there are some improvements of the Lagrangian interpolation in
the literature: \cite{schaback1992} addresses the problem of numerical
differentiation with application to computer-aided-design (the basic
assumption that we know the range of $f$ made in that paper makes it
unrealistic in control theory, but still deserves some attention),
\cite{criscuolo1990,criscuolo1991,criscuolo1992a,criscuolo1992b} propose an
{\it osculatory} alternative by assuming available some extra data points
which represent values of the higher derivatives of $f$ at the sampling
instants (this again does not directly fit with our system theory problem).

\smallskip\noindent
{\bf 2.2.2. Differentiating the spline interpolant.}
Excellent introductions to spline theory may be found in
\cite{greville1969b,nonweiler1984}; \S 3 of the latter book is particularly
recommendable for a motivated step by step introduction of splines in the
theory of interpolation.

A {\it natural spline} of {\it degree} $2l-1, \, l \in \setN, \, l \ge 1$
with {\it knots} $t_0, t_1, \cdots, t_N$ is a real valued function, $s$,
defined on $\setR$, and such that:
\item{(i)} $s$ has $2l-2$ first derivatives defined and continuous on
$\setR$,
\item{(ii)} In each interval $\left[ t_i,t_{i+1} \right], \, i = 0, 1,
\cdots, N-1$, $s$ is polynomial and with degree at most $2l-1$,
\item{(iii)} In each interval $\left] \infty,t_0 \right], \left[t_N,\infty
\right[$ $s$ is polynomial with degree at most $l-1$.

The set of natural splines of degree $2l -1$ with knots $t_0$, $t_1,
\cdots$, $t_N$ is classically denoted by $\CN_{2l-1}(t_0,$ $t_1, \ldots$,
$t_N)$.
It is clear that $\CN_{2l-1}(t_0, t_1, \ldots, t_N)$ contains the set
of real coefficient polynomials in $t$ with degree less than $l$.

The fine points of spline interpolation theory are the following
fundamental properties, see
\cite{schoenberg1964d,deboor1966,greville1969b}, for example, for more
properties and proofs.

\proclaim Interpolation property.
Let $f$ be in $\BE$.
For any natural integer $l$ such that $l \le N+1$ there is
one and only one natural spline $\hat s$ in $\CN_{2l-1}(t_0, t_1, \ldots
t_N)$ which interpolates $f$.
When $l = N+1$ then, of course, $\hat s$ is the Lagrange polynomial which is
the unique polynomial of degree less than $N+1$ interpolating $f$.
For $l \ge N+1$, there are infinitely many natural splines which
interpolate $f$.

\proclaim Smoothing properties.
Given $f$ in $\BE$, and $l \le N+1$ then
$$
\int_\BI \left( s^{(l)}(t) \right)^2 dt \le
\int_\BI \left( f^{(l)}(t) \right)^2 dt
$$
for any $s$ in $\CN_{2l-1}(t_0, t_1, \ldots, t_N)$, and equality holds
only if $f$ is the unique natural spline $\hat s$ interpolating $f$.

The reason why this last property is so-qualified {\it smoothing} is
explained very clearly in \S 3.4 of \cite{nonweiler1984}. 
It refers to {\it cubic splines} where the above integrals may be
interpreted in terms of energy.
Note that this minimum property may be generalized a bit farther as is done
in Sard's linear approximation theory~\cite{sard1967}, see also
\cite{greville1969b}, and
\cite{micchelli1976c,micchelli1976a,micchelli1976b}.
Roughly speaking, the generalization consists in substituting a general
linear functional for $f^{(l)}(t)$ in the above integrals.
We need the theory of Peano kernels, or more generally, of reproducing
kernels in Hilbert spaces in order to understand and use these
generalizations of spline interpolation.
Anyhow, the idea of smoothing
should not remain, instead, we have that of minimization.

Let us also ask if an ultimate generalization to the nonlinear case would
not be a definite solution to Problem~P, see \cite{micchelli1976b} for what
seems to be a step towards that idea, and also the references therein.

We postpone to the next version of this draft the summary of the results
on error bounds of spline numerical differentiation.
Let us just mention some papers containing partial results:
\cite{golomb1959,varma0000,birkhoff1964,secrest1965,sard1967,greville1969b,dolezal1982}.
We have not yet consulted \cite{deboor1978b} which is reportedly expect to
be an excellent source for spline interpolation.
This last paper contains an
innovative idea which consists in improving the spline interpolation by
iterating the basic process of spline interpolation as many times as we are
looking for a better precision. It is claimed that arbitrary precision may
thus be obtained.

Spline interpolation certainly provides better estimation, and constitutes
the most elaborate method of interpolation when one restricts oneself to
linear approximation and with polynomial base functions.

Beyond that method one has to enter the theory of {\it nonlinear
approximation} (see~\cite{rice1965}) with its simplest form, the {\it
rational approximation.}
But, then caution is recommended in the literature,
and, anyway, results are far less simple and clear.

The error bounds, say $\varepsilon_0$, $\varepsilon_1, \ldots$, on the
respective derivatives $f^{(0)}$, $f^{(1)}, \ldots$ of a function $f$,
no matter what the interpolation process is, are reasonably
expected to be such that $\varepsilon_0 <$ $\varepsilon_1 < \ldots$,
see~\cite{czipszer1958}.
That would mean that the lower the derivatives of $f$ to be estimated are the
easier we could afford high accuracy on the estimates of the derivatives.

Finally, let us note that we do yet not have at hand the references
\cite{newman0000,anderssen1972} which seem to be very
interesting\footnote\fnast{\sevenpoint It is remarkable that these
researchers were with engineering departments such as the IBM
Research Center at Yorktown Heights, NY; some other contributors on
numerical differentiation used to be with Ford Mo.~Co., Dearborn, MI.}.

\medskip\noindent
{\elevenbf 2.3. Numerical differentiation by regularization.}
As we said earlier the only reference on this approach that we yet have at
hand is the paper by Cullum~\cite{cullum1971}.
This paper mentions some other works by Russian numerical analysts.
Another bibliographic research is also due since there might be other
works which have been completed since 1971.

Let us just say here that this approach seems to be the most promising in
the standpoint of system theory.
Our problem~P is merely tackled in its general form.
Nevertheless, Cullum is concerned with the approximation of $f^{(1)}$ from
the samples of $f$.
Are her results generalizable to the approximation of finitely many
derivatives of $f$, and if yes, at what expense?
That is the question!

\bigskip\noindent
{\twelvebf 3. Illustrative examples.}

In order to better understand some of the issues involved in the use of  
numerical differentiation and interpolation in observer design, three academic
examples are currently being pursued.  The first example simply involves the estimation
of a time function and a few of its derivatives from noisy samples collected at
discrete instants of time.  The second example builds upon the first one
by showing that if the signal in question is generated by a known,
continuous-time linear model,
then the model's dynamcis can be incorporated into the estimation
process with advantage.  This will bring us into contact with more classical observer
theory.  The third example will illustrate that the technique is
also applicable to nonlinear system models, which is of course, the main
point.


In each case, polynomials are used as the interpolating
funtions, total least squares at the interpolating points is the 
metric to be minimized, and the discrete data will be gathered at a uniform rate.
All of  these simplifications are done so as to focus the discussion of other issues.

\medskip\noindent
{\elevenbf 3.1. Sum of two sinusoids.}

The purpose of this example is to illustrate the use of numerical
differentiation, alone and 
in combination with standard system theoretic notions,
to recover or estimate several derivatives of a signal in noise.  For numerical
differentiation, the key
parameters of interest are: $N$, the order of the interpolating polynomial;
$\Dt$, the time interval between data points;
$W * Dt$, the length in time of the moving window used for data interpolation
(note that $W + 1$ is the number of data points in the window);
$K$, the data node or ``knot'' within the moving window 
where the derivative is to be estimated; and $\s^2$, the ``intensity'' of the noise.

%The insight developed here will be transferred to an observation
%problem for a dynamical system.

Let an analog signal be given by the sum of two sinusoids of frequency 1 and
5 Hz., respectively:
\begin{equation}
y(t) = sin(2 \pi \w t) + sin(10 \pi \w t).
\label{sines}
\end{equation}
Thus $y$ could easily be generated by a dynamical system. 


Assume that the measured signal, $y^m$, is the sum of $y$ and $n$, where n is a
Gaussian white noise process, with variance $\s^2$.  It is further supposed that
$y^m$ is to be sampled at discrete instants of time,
$k \Dt$, $k = 0, 1, ...$.  In the absence of noise, the minimum sampling rate 
to reconstruct $y$ from sampled data would be the Nyquist rate, 10 Hz.; with
noise,
at least twice this rate is warranted, or 20 Hz..  As is standard pratice, an 
anti-aliasing filter will be used; its passband will be chosen as 20 Hz.  
It is then good practice to oversample the filtered signal, by a factor
of two in order to account for the non-ideal nature of the

anti-aliasing filter.   Discrete samples of $y + n$ will be collected therefore
at 40 Hz.,
after passing the signal through a fifth order Butterworth filter:
\begin{eqnarray}
y^m(t) &=& G_{BW}(s) (y(t) + n(t))
y^m_k & =& y^m(k \Dt),
\label{ysampled}
\end{eqnaray}
where $G_{BW}$ is the Butterworth, anti-aliasing filter.

Let the interpolating polynomial for the data $ \{ y_{k-W}^m, \ldots, y_k^m \}
be denoted by
\begin{eqaution}
\hat{y_k(t)} = a_0 + a_1 (t - t_{k-W}) + \ldots + a_N (t - t_{k-W})^N.
\label{interp_poly}
\end{eqaution}
The coefficents $\{ a_0, a_1, \ldots, \a_N \}$ are determined from the least
squares solution of 
\begin{eqnarray}
\left[ \begin{array}{cc} 1 & 0 & \ldots & 0 \\
1 & \Dt & \ldots & (\Dt)^N \\
\vdots & \vdots & \vdots & \vdots \\
1 & W |\Dt & \ldots & (W \Dt)^N 
\end{array} \right] \left[ \begin{array}{c} a_0 \\ a_1 \\ \vdots \\ a_N
\end{array} \right] & = & \left[ \begin{array}{c} y^m_{k-W} \\
y^m_{k-W+1} \\ \vdots \\ y^m_k
\end{array} \right].
\label{least_squares}
\end{eqnarray}
with respect to the Euclidean norm.  
The estimates of the derivatives of $y$ are determined by
\begin{equation}
\widehat{ \frac{d^j}{dt^j} y_k( \bar{t} )} := \frac{d^j}{dt^j} y_k(t)_{|t= \bar{t}};
\label{deriv_estimates}
\end{equation}
for simplicity of notation, this is written as 
$\frac{d^j}{dt^j} \hat{y}_k( \bar{t} )}$.


A number of numerical experiments or simulations
were conducted to ascertain the effects of the
parameters $N$,  $\D t$, $W$, $K$ and $\s^2$ on the ability to estimate the derivatives
of $y$.  A few of these are summarized here.  

{\bf Anna, this is your job!  (1) Discuss how the fitting was done.
(2) Discuss some of the various tradeoffs: order of
the
polynomial versus accuracy of the various derivatives and length of the window.
(3) Conclude with the ``best overall design'', maybe showing a few plots,
judiciously chosen.}

In order to make the error plots more meaningful, the errors for 
$(y, \ldots, frac{d^3}{dt^3})$ were normalized
by $(0.796, 19.3, 650, 22,700)$ respectively, which represent
$$lim_{T \rightarrow \infty} \frac{1}{T} \int_{0}^{T}{|\frac{d^j}{dt^j}y(t)| dt}$$.

The above developments described essentially an FIR (finite impulse response) 
filter, since the
estimates of the signal and its derivatives were completely determined by
the data points in the moving window of length $W$.  A kind of IIR (infinite
impulse response) filter can be achieved by introducing a simple modification
to the above; the modification is somewhat analogous to what is done in spline
approximations, but to a systems person, would be called a state.

The idea is to append to (ref{least_squares}), the estimates of $y$ and
its derivatives at node K from the previous window, that is, one appends
the relations
\begin{equation}
\frac{d^j}{dt^j} \hat{y}_k(t_{k-W+K-2}) =  \frac{d^j}{dt^j} \hat{y}_{k-1}(t_{k-W+K-2}).
\label{deriv_constraints}
\end{equation}
This is equivalent to
\begin{eqnarray}
\left[ \begin{array}{cc} 1 & 0 & \ldots & 0 \\
1 & (K-2) \Dt & ((K-2)\D t)^2 & \ldots & ((K-2) \D t)^j & \ldots & ((K-2) \Dt)^N \\
0 & 1 & 2(K-2) \D t & \ldots & j((K-2) \D t)^{j-1} & \ldots & N((K-2) \Dt)^{N-1} \\
\vdots & \vdots & \vdots & \vdots & \vdots & \vdots & \vdots \\
0 & 0 & \ldots & \ldots & j! & \ldots & N(N-1)\ldots(N-j+1)((K-2) \D t)^{N-j}
\end{array} \right] \left[ \begin{array}{c} a_0 \\ a_1 \\ \vdots \\ a_N
\end{array} \right] & = & \left[ \begin{array}{c} 
\hat{y}_{k-1}(t_{k-W+K-2}) \\
\frac{d}{dt} \hat{y}_{k-1}(t_{k-W+K-2}) \\ 
\vdots \\ 
\frac{d^j}{dt^j} \hat{y}_{k-1}(t_{k-W+K-2})
\end{array} \right].
\label{matrix_deriv_constraints}
\end{eqnarray}

The right-hand side of (\ref{deriv_constraints}) becomes a state which must be
initialized.  It has the effect of smoothing the fit from window 
to window; often, a smaller window can be used to achieve similar
noise attenuation.  
Weights can easily be added to trade-off errors in interpolating the data points
versus errors in interpolating the derivatives.  Figure \ref{??} shows the
effect of this on the estimation process for $N= ?$, $\D t = .025$, 
$W=?$ and $K=?$; the weight was chosen as ??.


\medskip\noindent
{\elevenbf 3.2. Including a model in the interpolation process.}

The goal of this example is to demonstrate two methods for incorporating
a dynamic model into the interpolation/numerical differentiation process.
This will bring us into even more direct contact with observer design because we
will tightly connect estimating the derivatives with estimating the states
of the model.

The signal (\ref{sines}) can obviously be generated by an initialized,
uncontrolled, linear model with four states.  Let 
\begin{eqnarray}
\frac{d}{dt} x & = &  A x \\
y & = & C x
\label{linear_model}
\end{eqnarray}
be such a model; it will be observable.  This means that $x$ can be recovered
from estimates of y and its derivatives through
\begin{equation} \left[ \begin{array}{c} 
y \\ \frac{d}{dt}y \\ \vdots \\ \frac{d^{\eta)}{dt^{\eta}}y
\end{array} \right] = \left[ \begin{array}{l} 
C \\ CA \\ \vdots \\ CA^{\eta}
\end{array} \right] x
\label{obs_mat}
\end{equation}
for $\eta$ large enough.  

The first way to include the model in the estimation process is now easily
explained.  Suppose that in general the model (\ref{linear_model}) has dimension
$n$, and for simplicity, assume that it has a single output.
As before, let
\begin{equation}
\left[ \begin{array}{c} 
\hat{y}_{k-1}(t_{k-W+K-2}) \\
\frac{d}{dt} \hat{y}_{k-1}(t_{k-W+K-2}) \\ 
\vdots \\ 
\frac{d^{n-1}}{dt^{n-1}} \hat{y}_{k-1}(t_{k-W+K-2})
\end{array} \right].
\end{equation}
be estimates of $y$ and its first $n-1$ derivatives.
From (\ref{obs_mat}) with $\eta = n-1$, determine $\hat{x}_{k-1}(t_{k-W+K-2})$. 
If the estimate of x is ``correct'', then it satisfies (\ref{obs_mat}) for
any $\eta$; if not, choosing $\eta \ge n$ introduces constraints whose
errors are to be minimized in order for x to be more closely compatible
with the dynamics (\ref{linear_model}).  Thus, in place of
(\ref{deriv_constraints}),
one augments (\ref{least_squares}) with
\begin{equation}
\frac{d^j}{dt^j} \hat{y}_k(t_{k-W+K-2}) = CA^j  \hat{x}_{k-1}(t_{k-W+K-2}),
\label{model_constraints}
\end{equation}
and procedes as indicated previously.  Note that taking $\eta = n-1$ reduces 
(\ref{model_constraints}) to (\ref{deriv_constraints}).

The method just outlined for determining $x$ from y and its derivatives is akin to
a ``dead-beat'' observer, though there is smoothing in the
interpolation of $y$, which is ``IIR'', as explained
earlier.  Some additional smoothing from the
model can be obtained, and constitutes the second method for introducing
the model into the interpolation/numerical differentiation process; moreover,
this technique, which is reminiscent of \cite{GM91, MG93}, will be
directly applicable to nonlinear systems.  

Let $A_d$ denote the discretization of the linear model, (\ref{linear_model}),
and let $\hat{x}_{k-1}(t_{k-W+K-2})$ be the previous estimate of $x$.  Update $x$
based on the latest estimate of $y$ and its derivatives through a damped Newton
method:
\begin{equation}
\hat{x}_{k}(t_{k-W+K-1} = A_d \hat{x}_{k-1}(t_{k-W+K-2}) + 
\epsilon \left[ \begin{array}{l} C \\ CA \\ \vdots \\ CA^{\eta} \end{array} 
\right]^{-1} ( \left[ \begin{array}{c} 
\hat{y}_{k}(t_{k-W+K-1}) \\
\frac{d}{dt} \hat{y}_{k}(t_{k-W+K-1}) \\ 
\vdots \\ 
\frac{d^{n-1}}{dt^{n-1}} \hat{y}_{k}(t_{k-W+K-1}) \end{array} \right] - 
\left[ \begin{array}{l} C \\ CA \\ \vdots \\ CA^{\eta} \end{array} 
\right] A_d \hat{x}_{k-1}(t_{k-W+K-2}) ),
\label{newton}
\end{equation}
where $0 < \epsilon <  1$; using $\epsilon = 1$ reduces it to a ``dead-beat''
estimator for $x$.  The damped Newton update can be used alone or with any of
the modifications to (\ref{least_squares}) discussed previously.

For simplicity in the numerical simulations that follow,
we chose the state in (\ref{linear_model}) as 
$$ x = \left[ \begin{array}{c}
y \\ \frac{d}{dt} y \\  \frac{d^2}{dt^2} y \\ \frac{d^3}{dt^3} y 
\end{array} \right]. $$  
It was properly initialized so that its output is precisely
given by (\ref{sines}).  The same noise sequence used in
Figure \ref{??} was added to the output.


Figure \ref{ ??} shows the results of .....


\medskip\noindent
{\elevenbf 3.3. Third order, nonlinear dynamical system.}

The goal of this example is to show that the method illustrated on the
linear example above also is applicable to nonlinear models.  For this
purpose, consider the following system which is a ``convex combination''
of an unstable linear dynamcis and a stable linear dynamics:
\begin{eqnarray}
\frac{d}{dt} \left[ \begin{array}{c} x_1\ \\ x-2 \\ x_3 \end{array} \right] 
& = & f(x) = 
\left[ \begin{array}{c} x_2 \\ x_3 \\ f_3(x)
\end{array} \right] \\
y & = & h(x) = x_1 ,\\
\mbox{where} & & \\
f_3(x) & = & \alpha(x) g_s(x) + (1 - \alpha(x)) g_u(x), \\
\alpha(x)& =& \beta  \frac{x'x}{1 + x'x} \\ 
g_s(x) & = & -54 x_1 - 36 x_2 - 9 x_3 \\
g_u(x) & = & 54 x_1 - 36 x_2 + 9 x_3.
\label{nonlin_model}
\end{eqnarray}
The parameter $\beta$ is nominally equal to 2, and in the final version of the
paper, will be be assumed unknown.  For now, it will be fixed at its nominal
value.

A state estimator was constructed using the interpolation/numerical
differentiation method of (\ref{least_squares}), with the model 
constraints (\ref{model_constraints}) replaced by
\begin{equation}
\frac{d^j}{dt^j} \hat{y}_k(t_{k-W+K-2}) = L_f^j h(\hat{x}_{k-1}(t_{k-W+K-2})),
\label{nonlin_model_constraints}
\end{equation}
and (\ref{newton}) replaced by
\begin{eqnarray}
\hat{x}_{k}^{-}(t_{k-W+K-1} &=& F_{\D t}(\hat{x}_{k-1}^{+}) \\
\hat{x}_{k}^{+}(t_{k-W+K-1} & = & \hat{x}_{k}^{-}(t_{k-W+K-1} + 
\epsilon \frac{\del}{del x}
\left[ \begin{array}{l} h \\ L_fh \\ L_f^2h \end{array} 
\right]^{-1}_{|\hat{x}_{k}^{-}(t_{k-W+K-1}} ( \left[ \begin{array}{c} 
\hat{y}_{k}(t_{k-W+K-1}) \\
\frac{d}{dt} \hat{y}_{k}(t_{k-W+K-1}) \\ 
\frac{d^{n-1}}{dt^{n-1}} \hat{y}_{k}(t_{k-W+K-1}) \end{array} \right] - 
\left[ \begin{array}{l} h \\ L_fh \\ L_f^2h \end{array} 
\right]_{\hat{x}_{k}^{-}(t_{k-W+K-1}) ).
\label{nonlin_newton}
\end{eqnarray}
In the above, $F_{\Dt}(x)$ is ..........


Figure \ref{?} shows the .........for $N = ??$ ............



\goodbreak\bigskip\noindent
{\twelvebf References}

{\sevenpoint
Though far from being exhaustive, the list of references is
intentionally prohibitive: all known potentially helpful piece of
literature is included.
Many of them have not been quoted in this first draft.
The starred ones are even not yet at our disposal.
The final version should consider only a few of
these references.
}

\bibliography{numdiff}
\bibliographystyle{plain}
\bye



